\documentclass[10pt]{article}
\usepackage[margin=0.60in]{geometry}
\usepackage{booktabs,enumitem,fancyhdr,graphicx,listings,tabularx,titlesec,xcolor}
\usepackage[hidelinks]{hyperref}
\graphicspath{{docs/lessons/assets/}{assets/}}
\hypersetup{
 pdftitle={ADK Lesson 067: Normalize a Copied Inertial Record},
 pdfauthor={Kevin Shanahan},
 pdfsubject={Stable typed records and canonical fixed images from copied inertial evidence},
 pdflang={en-US}}
\pagestyle{fancy}\fancyhf{}
\lhead{\textsc{ADK Lesson 067}}\rhead{Copied Inertial Records}
\cfoot{\thepage}\setlength{\headheight}{14pt}
\setlength{\parindent}{0pt}\setlength{\parskip}{3pt}
\setlist{nosep,leftmargin=1.45em}
\titleformat{\section}{\Large\bfseries}{\thesection}{0.5em}{}
\newcommand{\lessonbox}[1]{\begin{center}\fbox{\parbox{0.92\linewidth}{#1}}\end{center}}
\newcommand{\blankline}{\rule{0.76in}{0.4pt}}
\newcommand{\longblank}{\rule{1.35in}{0.4pt}}
\lstset{
 basicstyle=\ttfamily\footnotesize,
 columns=fullflexible,
 breaklines=true,
 frame=single,
 rulecolor=\color{black},
 showstringspaces=false}

\begin{document}
\small
\begin{center}
 {\Huge\bfseries Normalize a Copied Inertial Record}\\[3pt]
 {\Large Make evidence stable before asking what it means}\\[7pt]
 \fbox{\textbf{E0 SYNTHETIC REPLAY --- NO SENSOR, BUS, PIN, OR CALIBRATION CLAIM}}
\end{center}
\begin{tabularx}{\linewidth}{@{}>{\bfseries}l X >{\bfseries}l X@{}}
Lesson & 067 --- component & Time & 150--210 minutes\\
Prerequisite & Lesson 043 copied inertial samples & Board & compile-only Mega replay\\
Input & one copied sample & Observation & named record and 64-byte image\\
\end{tabularx}

\lessonbox{\textbf{Responsibility.}
\texttt{InertialRecordNormalizer} copies one complete Lesson 043 sample into
an explicit stable record. \texttt{InertialRecordCodec} gives that record one
canonical 64-byte representation. Neither object reads a sensor, rotates an
axis, removes bias, qualifies a source, discovers hardware, or writes storage.}

\section{Draw the evidence boundary first}
A measurement is useful only when its identity and conditions travel with its
numbers. The copied sample already names its source, declared ranges,
calibration revision, observation time, sequence, readiness, saturation, and
producer status. Lesson 067 preserves that whole witness. It does not rescue a
fault by borrowing vectors from the preceding healthy sample.

\begin{center}
% ADK visual: pencil
\includegraphics[width=0.96\linewidth,height=2.10in,keepaspectratio]
 {067-copied-boundary-pencil.png}

\small\textit{Alternative text: a copied sample card crosses a normalization
desk into a stable typed record and fixed image; a physical sensor, bus, pins,
and calibration remain behind a locked E1 gate with no connecting wire.}
\end{center}

\begin{tabularx}{\linewidth}{@{}X X@{}}
\toprule
E0 establishes & E0 does not establish\\
\midrule
field-for-field copied provenance & genuine MPU6050 or QMI8658 identity\\
stable revisions and closed states & physical scale, axes, bias, or accuracy\\
canonical bytes and corruption checks & I2C behavior, voltage, timing, or pins\\
deterministic replay & qualification, mounting, or interchangeability\\
\bottomrule
\end{tabularx}

\section{Predict before normalizing}
For a healthy, ready sample, predict the derived record state:
\longblank. If the sample is healthy but not ready: \longblank. Write the
fields that must remain visible when producer status reports a fault:
\hrulefill.

\newpage
\section{Derive one state without discarding evidence}
The record has three derived states. \texttt{Recorded} means producer status
is OK and data is ready. \texttt{NotReady} means producer status is OK and
data is not ready. Any non-OK producer status becomes \texttt{SourceFault}.
Saturation remains a separate field; it is never rewritten as producer status.

\begin{center}
% ADK visual: pencil

\includegraphics[width=0.94\linewidth,height=2.12in,keepaspectratio]
 {067-state-triad-pencil.png}

\small\textit{Alternative text: three graphite cards show Recorded for
healthy ready evidence, NotReady for healthy non-ready evidence, and
SourceFault for producer failure; saturation remains visible beside every
state.}
\end{center}

\begin{tabularx}{\linewidth}{@{}l l l X@{}}
\toprule
Producer status & Ready & Derived state & Numeric vectors\\
\midrule
OK & yes & \texttt{Recorded} & copied exactly\\
OK & no & \texttt{NotReady} & copied exactly\\
non-OK & either & \texttt{SourceFault} & copied exactly, never cached\\
\bottomrule
\end{tabularx}

This separation matters. ``Not ready'' is a valid report about an identified
source; it is not the same as a transport failure. Saturation is valid
evidence that a declared measurement range was exceeded; it must not be
averaged away or mistaken for an ordinary maximum.

\section{Keep revisions nonzero and explicit}
\texttt{schemaRevision} describes the typed record contract.
\texttt{normalizationRevision} describes this field-preserving transformation.
Both must be nonzero. Source identity retains kind, model, source ID, source
configuration revision, and calibration revision. A future lesson may compare
or qualify records, but Lesson 067 merely ensures that comparison has complete
inputs.

Circle each change that should produce a different canonical record: source
ID; calibration revision; acceleration range; readiness; sequence; status.
Explain why ``the vectors are equal'' is not enough: \hrulefill.

\newpage
\section{Lay out the canonical image deliberately}
The codec writes exactly 64 bytes in little-endian order. It is a designed
wire/storage-neutral image, not the in-memory layout of a C++ object.
Magic, version, and encoded length establish framing. Reserved bytes are zero.
A CRC-16 with polynomial \(0x1021\), initial value \(0xffff\), and final XOR
\(0x0000\) covers bytes 0 through 61.

\begin{center}
% ADK visual: pencil
\includegraphics[width=0.98\linewidth,height=2.30in,keepaspectratio]
 {067-record-layout-pencil.png}

\small\textit{Alternative text: a hand-drawn 64-cell notebook strip groups
framing, revisions, identity, ranges, two three-axis vectors, observation time,
sequence, flags, status, and the final two-byte CRC seal.}
\end{center}

\begin{tabularx}{\linewidth}{@{}r l X@{}}
\toprule
Bytes & Group & Required interpretation\\
\midrule
0--5 & framing & magic, image version, encoded length\\
6--9 & contract & schema and normalization revisions\\
10--19 & identity & state, complete source, calibration, reserved zeroes\\
20--27 & ranges & declared acceleration and angular-rate ranges\\
28--51 & vectors & signed x/y/z values in Lesson 043 units\\
52--59 & ordering & supplied modular time and sequence\\
60--61 & quality & ready/saturation flags and producer status\\
62--63 & integrity & little-endian CRC-16 over bytes 0--61\\
\bottomrule
\end{tabularx}

Mark the first byte range you would inspect for each symptom: wrong model
\blankline; changed vector \blankline; bad CRC \blankline; unknown flag bit
\blankline.

\section{Do not confuse a checksum with trust}
The CRC detects accidental change according to this format. It does not
authenticate a producer, establish physical calibration, make volatile bytes
durable, or prove that a sensor supplied them. Full typed fields remain the
authoritative evidence; the CRC is an integrity check over their canonical
encoding.

\newpage
\section{Stage the entire result before publishing}
\texttt{normalize()} and codec operations are transactional. They build and
validate a complete candidate first. An invalid configuration, invalid closed
enum, short output span, inconsistent state, or corrupt image leaves the
caller's previously published output byte-for-byte unchanged.

\begin{center}
% ADK visual: pencil
\includegraphics[width=0.94\linewidth,height=2.18in,keepaspectratio]
 {067-atomic-desk-pencil.png}

\small\textit{Alternative text: an old record and image wait beside a staging
desk; only a fully checked candidate receives the pass mark and replaces both,
while a failed candidate leaves the old output and its canary untouched.}
\end{center}

\begin{enumerate}
 \item Copy every field into a candidate typed record.
 \item Derive state from producer status and readiness.
 \item Validate revisions, closed enums, flags, and semantic combinations.
 \item Encode all fields at locked offsets and calculate the CRC.
 \item Commit the complete candidate only after every check passes.
\end{enumerate}

\lessonbox{\textbf{Atomicity experiment.}
Fill an output record and surrounding canaries with a recognizable pattern.
Try each malformed sample or image. A correct implementation reports the
failure and preserves the record and both canaries exactly.}

\section{Compare fields, not padding}
Semantic record equality compares named fields. \texttt{memcmp} can observe
padding, object representation, or compiler choices that are not part of the
contract. Canonical image equality compares all 64 encoded bytes because that
format deliberately assigns meaning to every byte, including reserved zeroes.

\newpage
\section{Diagnose rejection precisely}
Decode first checks length, then framing, integrity, and semantic values before
publishing. The typed result helps a learner locate the failed boundary without
turning diagnostic text into library state.

\begin{center}
% ADK visual: pencil
\includegraphics[width=0.97\linewidth,height=2.16in,keepaspectratio]
 {067-corruption-lab-pencil.png}

\small\textit{Alternative text: four scratched record strips fall into
separate diagnosis trays for BadLength, BadFraming, BadIntegrity, and
BadSemanticValue; rejected strips cannot alter the prior typed output.}
\end{center}

\begin{tabularx}{\linewidth}{@{}l X X@{}}
\toprule
Validity & Example & What remains unchanged\\
\midrule
\texttt{BadLength} & 63 or 65 input bytes & typed output\\
\texttt{BadFraming} & magic, version, length byte, flags, reserved byte & typed output\\
\texttt{BadIntegrity} & payload changed without matching CRC & typed output\\
\texttt{BadSemanticValue} & unknown enum or contradictory state & typed output\\
\bottomrule
\end{tabularx}

Prediction: flip one vector bit and do not update the CRC. The result should be
\longblank. Now recompute the CRC while leaving an invalid state/ready
combination. The result should be \longblank. This shows why integrity and
meaning are independent checks.

\section{Use status and validity for different questions}
\texttt{Status} reports operation failures such as invalid configuration or
insufficient destination capacity. \texttt{InertialRecordValidity} classifies
an attempted decode. Neither type claims sensor health; producer health remains
inside the record.

\newpage
\section{Replay without touching hardware}
The Mega example follows acquire, configure, start, then observe, decide,
actuate. It acquires one synthetic copied sample, configures two nonzero
revisions, normalizes the sample, encodes the record, and exposes named
volatile record and image cells. Here ``actuate'' means publish inert values;
no endpoint is invoked.

\begin{center}
% ADK visual: pencil
\includegraphics[width=0.95\linewidth,height=2.02in,keepaspectratio]
 {067-replay-ledger-pencil.png}

\small\textit{Alternative text: two identical copied samples pass through two
normalization runs and produce field-equal typed records and byte-equal
64-byte strips.}
\end{center}

\begin{enumerate}
 \item \textbf{Predict:} write the expected record state, identity, first
       vector component, sequence, flags, and image length.
 \item \textbf{Observe:} compile for Mega 2560 and inspect the named volatile
       record and image cells with a debugger, simulator, or disassembly.
       Serial may support this observation but is not required.
 \item \textbf{Interpret:} compare every named field and all 64 bytes across
       two runs. Explain any mismatch before changing the fixture.
\end{enumerate}

The named copied record and image are the E0 non-Serial observation path.
They establish deterministic software publication separately from resource
acquisition and electrical safe state. E0 owns no resource or endpoint, so
those physical checks remain open rather than silently passing.

\section{Test the boundaries, not just the example}
\begin{itemize}
 \item Exercise every supported source, state, readiness, saturation, range,
       and producer-status class.
 \item Preserve negative, zero, positive, and extreme accepted components.
 \item Reject zero revisions, unknown enums, unknown flags, nonzero reserved
       bytes, inconsistent semantics, and every length error.
 \item Use byte-exact golden images, round trips, individual field corruption,
       CRC corruption, and output canaries.
 \item Prove that adding an unrelated later library component does not change
       Lesson 067 records or its resource evidence.
\end{itemize}

\newpage
\section{Draw the line around the next lesson}
Lesson 068 may qualify exactly one explicitly configured source and map its
copied vectors through one proper axis rotation. It must consume the complete
Lesson 067 record. Lesson 067 itself does no axis mapping, stationary-window
check, source selection, voting, failover, or pitch/roll presentation.

There is no pin-by-pin connection list, wiring diagram, or formal schematic at
E0 because there is no circuit. A future E1 adapter first needs an exact chip
and carrier revision, primary data sheet and register map, identity and reset
behavior, address straps, byte order, ranges and scaling, data-ready behavior,
mounting axes, regulator and pull-ups, logic levels, current, bus ownership,
rollback, safe state, and bench acceptance. Evidence for one device family
cannot close another family's card.

\section{Troubleshooting matrix}
\begin{tabularx}{\linewidth}{@{}l X X@{}}
\toprule
Symptom & Inspect first & Interpretation\\
\midrule
\texttt{NotReady} & producer status and ready flag & valid copied evidence; no new ready sample\\
\texttt{SourceFault} & complete producer status & source reported failure; vectors are not rescued\\
saturation set & declared range and saturation field & quality warning remains independent of state\\
\texttt{BadFraming} & magic/version/length/flags/reserved & image cannot enter the semantic boundary\\
\texttt{BadIntegrity} & bytes 0--61 and stored CRC & canonical bytes changed\\
\texttt{BadSemanticValue} & closed enums and state consistency & intact bytes encode an invalid meaning\\
\bottomrule
\end{tabularx}

\section{Exercises}
\begin{enumerate}
 \item Explain why a fault record must preserve its current vectors rather
       than repeat the last healthy vectors.
 \item Hand-encode one negative 32-bit vector component in little-endian order.
 \item List three edits that change the CRC but not the image length.
 \item Design a changed-duplicate test in which time and sequence are equal but
       one identity field differs.
 \item Explain why an intact CRC is neither authentication nor physical
       qualification.
\end{enumerate}

\newpage
\section{Open acceptance record}
\begin{tabularx}{\linewidth}{@{}>{\bfseries}l X@{}}
Host evidence & Pending integration: strict host tests, no-exception tests,
standalone header, golden images, corruption matrix, atomicity, and canaries.\\
Mega evidence & Pending integration: compile the canonical synthetic replay
and record exact flash and static SRAM.\\
Memory evidence & Planning limits are a 4-byte normalizer and an exact 64-byte
canonical image; record measured object and synchronous-stack evidence before
publication.\\
Document evidence & Build twice; check metadata, embedded fonts, extraction,
monochrome output, links, classifications, and rendered pencil legibility.\\
Physical evidence & Open E1: no sensor, bus, pin, mounting, range, calibration,
data-ready, safe-state, or bench observation has been established.\\
\end{tabularx}

\vspace{8pt}
\begin{tabularx}{\linewidth}{@{}l X@{}}
\toprule
Acceptance item & Record\\
\midrule
Commit tested & \hrulefill\\[10pt]
Toolchain and Mega core & \hrulefill\\[10pt]
Host and sanitizer result & \hrulefill\\[10pt]
Flash / static SRAM / stack & \hrulefill\\[10pt]
Record object / normalizer / image sizes & \hrulefill\\[10pt]
Golden-image and corruption result & \hrulefill\\[10pt]
PDF visual and extraction review & \hrulefill\\[10pt]
Open E1 owner and next action & \hrulefill\\
\bottomrule
\end{tabularx}

\lessonbox{\textbf{Stop condition.}
Do not call this hardware verified. The completed E0 exercise proves a
deterministic record contract over copied values. Exact powered acquisition,
physical axes, calibration, observation, and safe state remain separate work.}

\vfill
\begin{center}
\textbf{Next: Lesson 068 --- Qualify one configured inertial record source.}
\end{center}
\end{document}
